\documentclass[12pt]{article}
\usepackage{amsmath}
\usepackage{amssymb}
\usepackage{geometry}
\usepackage{setspace}
\geometry{margin=1in}
\setstretch{1.2}

\title{Lecture Scribe \\ Joint Random Variables and Joint Distributions}
\author{Course: CSE400 -- Fundamentals of Probability in Computing}
\date{}

\begin{document}

\maketitle

\section{Introduction: Why Study Joint Random Variables?}

In science and real life, we are often interested in two or more random variables at the same time. Studying these variables together allows us to compute probabilities involving both variables and understand the relationship between them.

\subsection*{Examples of Situations Involving Two Variables}

\begin{itemize}
\item Height and weight of giraffes
\item IQ and birthweight of children
\item Frequency of exercise and rate of heart disease
\item Air pollution level and respiratory illness rate
\item Number of Facebook friends and age of members
\end{itemize}

These examples demonstrate that many real-world problems involve multiple variables whose behaviors are related.

\subsection*{Example 1: Smart Transportation}

In a smart traffic system two random variables are defined:

\[
X = \text{Vehicle Speed}
\]

\[
Y = \text{Traffic Density}
\]

Question raised:

\textit{Can speed be high when traffic density is also high?}

This question highlights that the variables may be correlated random variables, meaning their values may influence each other.

Applications include:

\begin{itemize}
\item Adaptive traffic signals
\item Congestion prediction
\item Autonomous vehicle routing
\end{itemize}

\subsection*{Example 2: Wireless Network Performance}

Two random variables are defined:

\[
X = \text{Signal-to-Noise Ratio (SNR)}
\]

\[
Y = \text{Data Throughput}
\]

Throughput depends on SNR, but both variables vary due to:

\begin{itemize}
\item Fading
\item Interference
\item Congestion
\end{itemize}

Question raised:

\textit{If SNR is high, will throughput always be high?}

Throughput may not always be high due to network congestion, showing that the variables have conditional relationships.

Applications include:

\begin{itemize}
\item 5G scheduling
\item Adaptive modulation
\item Network planning
\end{itemize}

\subsection*{Additional Examples}

\textbf{Weather Prediction}

\[
X = \text{Rainfall}
\]

\[
Y = \text{Temperature}
\]

\textbf{Drone Delivery}

\[
X = \text{Wind Speed}
\]

\[
Y = \text{Delivery Time}
\]

\textbf{Rolling Two Dice}

Let \(X\) and \(Y\) represent the values obtained when two dice are rolled.

Example questions:

\[
X + Y = 7
\]

\[
X > Y
\]

These require the use of joint probability mass functions.

\subsection*{Importance of Joint Distributions}

When random variables are considered together, they have a \textbf{joint distribution}. The joint distribution allows us to:

\begin{itemize}
\item Compute probabilities involving both variables
\item Understand the relationship between the variables
\end{itemize}

When variables are independent, analysis becomes simpler. When they are not independent, dependence can be measured using:

\begin{itemize}
\item Covariance
\item Correlation
\end{itemize}

\section{Discrete Case -- Joint Probability Mass Function (Joint PMF)}

\subsection*{Setup}

Suppose \(X\) and \(Y\) are discrete random variables.

Possible values:

\[
X = \{x_1, x_2, \dots , x_n\}
\]

\[
Y = \{y_1, y_2, \dots , y_m\}
\]

The ordered pair \((X,Y)\) takes values in the product set:

\[
\{(x_1,y_1),(x_1,y_2),\dots,(x_n,y_m)\}
\]

Thus every possible combination of values of \(X\) and \(Y\) represents a possible outcome.

\subsection*{Definition: Joint PMF}

The joint probability mass function is defined as:

\[
p(x_i, y_j) = P(X = x_i, Y = y_j)
\]

This gives the probability that:

\begin{itemize}
\item \(X\) takes value \(x_i\)
\item \(Y\) takes value \(y_j\)
\end{itemize}

simultaneously.

\section{Joint PMF Table Representation}

The joint PMF can be organized in a table format.

\begin{itemize}
\item Rows represent values of \(X\)
\item Columns represent values of \(Y\)
\item Each cell contains \(p(x_i,y_j)\)
\end{itemize}

\subsection*{Key Properties of Joint PMF}

\textbf{Probability Bounds}

\[
0 \le p(x_i, y_j) \le 1
\]

\textbf{Total Probability Condition}

\[
\sum_{i=1}^{n}\sum_{j=1}^{m} p(x_i, y_j) = 1
\]

This means the sum of probabilities of all possible outcomes must equal 1.

\section{Example: Rolling Two Dice}

\subsection*{Experiment}

Roll two fair dice.

Define random variables:

\[
X = \text{value on the first die}
\]

\[
T = \text{total (sum) of both dice}
\]

Each pair outcome \((i,j)\) has probability:

\[
\frac{1}{36}
\]

Possible values:

\[
X = 1,2,3,4,5,6
\]

\[
T = 2,3,4,\dots,12
\]

Some combinations are possible with probability \(1/36\), while impossible combinations have probability \(0\).

\subsection*{Observations}

\begin{itemize}
\item Entries in the table are either \(0\) or \(1/36\)
\item \(X\) and \(T\) are not independent
\end{itemize}

Thus the value of \(T\) depends on the value of \(X\).

\section{Continuous Case -- Joint Probability Density Function (Joint PDF)}

The continuous case is analogous to the discrete case.

\begin{itemize}
\item Discrete values $\rightarrow$ Continuous intervals
\item Joint PMF $\rightarrow$ Joint PDF
\item Summations $\rightarrow$ Double integrals
\end{itemize}

\subsection*{Domain of Variables}

If

\[
X \in [a,b]
\]

\[
Y \in [c,d]
\]

then the pair \((X,Y)\) lies in the region

\[
[a,b] \times [c,d]
\]

\subsection*{Definition: Joint PDF}

A function \(f(x,y)\) is called a joint probability density function if

\[
P((X,Y) \text{ lies in a small rectangle } dxdy) \approx f(x,y)dxdy
\]

\subsection*{Key Properties}

\textbf{Non-negativity}

\[
f(x,y) \ge 0
\]

\textbf{Total Probability}

\[
\int_{c}^{d}\int_{a}^{b} f(x,y)\,dx\,dy = 1
\]

Note: the joint PDF may take values greater than 1 because it represents density.

\section{Continuous Joint PDF Example}

Given

\[
f(x,y) = 4xy
\]

with

\[
0 \le x \le 1
\]

\[
0 \le y \le 1
\]

\subsection*{Step 1: Check Validity}

Compute

\[
\int_0^1 \int_0^1 4xy\,dx\,dy
\]

First integrate with respect to \(x\):

\[
\int_0^1 4xy\,dx = 2y
\]

Then integrate with respect to \(y\):

\[
\int_0^1 2y\,dy = y^2 \Big|_0^1 = 1
\]

Thus \(f(x,y)\) is a valid joint PDF.

\subsection*{Step 2: Compute Probability}

Event:

\[
A = \{X < 0.5, Y > 0.5\}
\]

Probability:

\[
P(A) = \int_0^{0.5} \int_{0.5}^{1} 4xy\,dy\,dx
\]

After evaluation:

\[
P(A) = \frac{3}{16}
\]

\section{Joint Cumulative Distribution Function (Joint CDF)}

For jointly distributed random variables \(X\) and \(Y\),

\[
X \le x, \quad Y \le y
\]

represents the event:

\[
\{X \le x \text{ and } Y \le y\}
\]

\subsection*{Definition}

\[
F(x,y) = P(X \le x, Y \le y)
\]

This gives the probability that both random variables are less than or equal to specified values.

\section{Joint CDF -- Continuous Case}

If \(X,Y\) are continuous with density \(f(x,y)\):

\[
F(x,y) = \int_c^y \int_a^x f(u,v)\,du\,dv
\]

\subsection*{Recovering the Joint PDF}

\[
f(x,y) = \frac{\partial^2}{\partial x \partial y}F(x,y)
\]

\section{Joint CDF -- Discrete Case}

If \(X,Y\) are discrete with joint PMF \(p(x_i,y_j)\),

\[
F(x,y) = \sum_{x_i \le x}\sum_{y_j \le y} p(x_i,y_j)
\]

Interpretation:

Add all probabilities in the table with:

\begin{itemize}
\item rows \(x_i \le x\)
\item columns \(y_j \le y\)
\end{itemize}

\section{Properties of the Joint CDF}

\subsection*{Property 1: Non-Decreasing}

\(F(x,y)\) is non-decreasing. If \(x\) or \(y\) increases, \(F(x,y)\) either stays constant or increases.

\subsection*{Property 2: Lower-Left Corner}

\[
\lim_{(x,y)\to(-\infty,-\infty)} F(x,y) = 0
\]

\subsection*{Property 3: Upper-Right Corner}

\[
\lim_{(x,y)\to(\infty,\infty)} F(x,y) = 1
\]

\end{document}